\section{Validación Preliminar del Sistema}

Para asegurar la conformidad del producto con las especificaciones y requerimientos finales, se ha llevado a cabo una validación preliminar en un entorno que simula las condiciones operativas reales. Esta fase de pruebas es crucial para identificar posibles cuellos de botella y garantizar que el sistema de detección de billetes funcione de manera robusta y eficiente.

\subsection{Pruebas de Estrés y Rendimiento}

Las pruebas de estrés se diseñaron para evaluar el comportamiento del sistema bajo cargas de trabajo intensivas, midiendo los tiempos de inferencia y la estabilidad de la aplicación. Se ejecutaron pruebas variando la cantidad de imágenes procesadas simultáneamente para simular diferentes escenarios de uso real.

A continuación, se presenta una tabla que resume los resultados obtenidos durante las pruebas de rendimiento.

\begin{table}[H]
\centering
\caption{Resultados de Pruebas de Rendimiento del Detector de Billetes}
\label{tab:rendimiento_deteccion}
\begin{tabular}{cccc}
\toprule
\textbf{Imágenes} & \textbf{Inferencia (ms)} & \textbf{Threshold} & \textbf{Núcleos} \\
\midrule
20 & 290 ms & 0.8 & 2 \\
5  & 291 ms & 0.8 & 2 \\
1  & 314 ms & 0.8 & 2 \\
20 & 319 ms & 0.8 & 2 \\
5  & 336 ms & 0.8 & 2 \\
1  & 354 ms & 0.8 & 2 \\
10 & 379 ms & 0.8 & 2 \\
\bottomrule
\end{tabular}
\end{table}

Como se puede observar en la Tabla \ref{tab:rendimiento_deteccion}, los tiempos de inferencia se mantienen dentro de un rango aceptable, incluso al procesar un mayor número de imágenes. Esto demuestra la eficiencia del modelo y su capacidad para operar en condiciones de alta demanda.

\subsection{Pruebas de Estrés Adicionales Propuestas}

Para complementar la validación del sistema, se proponen las siguientes pruebas de estrés adicionales:

\begin{itemize}
    \item \textbf{Detección con Imágenes de Diferentes Resoluciones:} Evaluar cómo el tiempo de inferencia y la precisión del modelo se ven afectados al procesar imágenes con distintas resoluciones (e.g., 480p, 720p, 1080p). Esto permitiría determinar los límites operativos del sistema en diferentes dispositivos.
    \item \textbf{Pruebas de Detección bajo Diferentes Condiciones de Iluminación:} Simular condiciones de poca luz, sobreexposición y sombras para verificar la robustez del detector de billetes en escenarios de uso real.
    \item \textbf{Pruebas de Oclusión Parcial:} Analizar el comportamiento del modelo cuando los billetes están parcialmente cubiertos por otros objetos, para determinar el umbral de detección en condiciones no ideales.
    \item \textbf{Pruebas de Resistencia a Largo Plazo:} Ejecutar la aplicación de forma continua durante un período prolongado (e.g., varias horas) para detectar posibles fugas de memoria o degradación del rendimiento a lo largo del tiempo.
\end{itemize}

La implementación de estas pruebas permitiría obtener una visión más completa del rendimiento y la fiabilidad del sistema en un amplio rango de condiciones operativas.